\documentclass[UTF8,a4paper,11pt,oneside,openany]{ctexbook}
\usepackage{amsmath,amssymb,mathtools}
\usepackage{geometry}
\geometry{margin=23mm}
\usepackage[hidelinks]{hyperref}
\allowdisplaybreaks
\setlength{\parindent}{0pt}
\setlength{\parskip}{0.6em}
\begin{document}
\title{大学物理高难综合题（第33批候选）}
\author{Ruoshui Physics Bank}
\date{2026年10月}
\maketitle
\tableofcontents
\chapter*{重要质量说明}
这是一份审核分支候选资料；自动模型检查不等于独立专家逐式复核。
在跨批查重、人工复核及双册编译均完成前，严禁标记为正式主库验收。

\chapter{力学}
\section*{M220\quad Lagrange重对称陀螺的循环积分与睡眠顶稳定性}
\textbf{题目}\par
质量$m$的重对称陀螺定点支承，质心距支点$l$，惯量$I_1=I_2=I$、$I_3$；对称轴相对竖直向上角为$\theta$，Euler角$(\phi,\theta,\psi)$，重力势能$U=mgl\cos\theta$。
\par\smallskip\noindent (1) 推出Lagrangian及两个循环动量。
\par\smallskip\noindent (2) 在固定动量下化为一自由度的倾角有效势和能量方程。
\par\smallskip\noindent (3) 求匀速规则进动的平衡条件及相应进动角速度分支。
\par\smallskip\noindent (4) 讨论直立快速自转（睡眠顶）线性稳定临界角动量并检验重力零极限。
\par\medskip\hrule\medskip
\section*{M221\quad 旋转各向异性二维谐势的陀螺稳定区与负能模}
\textbf{题目}\par
质量$m$的质点限制在旋转平面内，实验室中一各向异性二次势以角速率$\Omega$绕竖直轴匀速转动；在旋转坐标中$U=\frac m2(\kappa_x x^2+\kappa_y y^2)$，其中$0<\kappa_x<\kappa_y$。
\par\smallskip\noindent (1) 写出旋转参考系的线性动力方程。
\par\smallskip\noindent (2) 求双圆偏振混合本征模对应的$\omega^2$方程。
\par\smallskip\noindent (3) 精确求线性谱稳定的转速区间，并说明临界转速的特殊性。
\par\smallskip\noindent (4) 计算Jacobi型守恒能量并说明高速稳定区为何不一定有正定能量。
\par\medskip\hrule\medskip
\section*{M222\quad 相对论变质量火箭的四动量平衡与双曲正切火箭方程}
\textbf{题目}\par
质量随燃料喷出而减小的理想直线火箭，初始静质量$M_0$、终质量$M_f$，喷气在火箭瞬时静止系中的后向速度大小恒定$u<c$；忽略外力和辐射。
\par\smallskip\noindent (1) 用瞬时共动系能量和动量守恒推导微分形式的速度增量。
\par\smallskip\noindent (2) 将速度复合写成快速度，积分得到火箭末速度。
\par\smallskip\noindent (3) 给出低速喷气与光子火箭极限。
\par\smallskip\noindent (4) 给定目标$v_f$反求燃料质量比，并检查燃料比趋一的极限。
\par\medskip\hrule\medskip
\section*{M223\quad 平面限制性三体问题的三角Lagrange点与Routh稳定阈值}
\textbf{题目}\par
规范化圆形限制性三体问题中，主星质量$1-\mu$和伴星质量$\mu$（$0<\mu\le1/2$）在旋转系固定为$(-\mu,0)$、$(1-\mu,0)$，旋转角速度为1；无质量粒子受引力及惯性力作用。
\par\smallskip\noindent (1) 建立有效势$\Omega(x,y)$及运动方程。
\par\smallskip\noindent (2) 验证三角平衡点$L_{4,5}$坐标。
\par\smallskip\noindent (3) 推导三角点二维线性扰动特征多项式。
\par\smallskip\noindent (4) 求谱稳定的临界质量比并解释等号共振极限。
\par\medskip\hrule\medskip
\section*{M224\quad 近圆轨道编队的Hill–Clohessy–Wiltshire方程与漂移条件}
\textbf{题目}\par
两小卫星围绕引力参数$\mu_G$的中心天体，主星半径$r_0$的圆轨道角速度$n=\sqrt{\mu_G/r_0^3}$。设从星相对主星的位置为$(x,y,z)$，分别指向径向、顺轨和轨道法向，且相对距离远小于$r_0$。
\par\smallskip\noindent (1) 从旋转系相对引力梯度推导Hill线性化方程。
\par\smallskip\noindent (2) 给出平面运动的完整通解与法向振动解。
\par\smallskip\noindent (3) 求顺轨线性漂移被抑制的初值条件。
\par\smallskip\noindent (4) 写出与方程一致的守恒二次量并说明线性模型的范围。
\par\medskip\hrule\medskip
\section*{M225\quad 恒定固有加速度的Rindler双曲线、视界与Doppler频移}
\textbf{题目}\par
一艘理想航天器在一维闵可夫斯基时空中从实验室静止出发，之后沿$+x$方向保持固有加速度$a>0$常量，初始$(t,x)=(0,0)$，光速$c$。
\par\smallskip\noindent (1) 以固有时$\tau$求四速度、实验室坐标和速度。
\par\smallskip\noindent (2) 验证世界线为双曲线，且固有加速度确为常量。
\par\smallskip\noindent (3) 讨论从实验室原点稍后发出的同向光信号能否追上火箭，求临界发射时刻。
\par\smallskip\noindent (4) 推导从后方同向传播光信号在火箭上的瞬时Doppler频移。
\par\medskip\hrule\medskip
\section*{M226\quad 引力双曲遭遇的散射角与行星引力弹弓能量交换}
\textbf{题目}\par
测试粒子以相对行星的无穷远速度$v_\infty>0$、碰撞参数$b>0$接近引力参数$\mu_p=GM_p$的行星，忽略行星以外引力在遭遇区的影响；行星在恒星系中的速度为常矢量$\boldsymbol V_p$。
\par\smallskip\noindent (1) 由二体比能量与比角动量求双曲离心率和近心距离。
\par\smallskip\noindent (2) 推导无穷远入、出速度方向的偏转角。
\par\smallskip\noindent (3) 在恒星系中推导引力弹弓的动能变化公式及其绝对上界。
\par\smallskip\noindent (4) 检验弱引力和正碰近极限，并说明拼接圆锥轨道的适用条件。
\par\medskip\hrule\medskip
\section*{M227\quad 引力波长波极限下的测地偏离与Michelson响应}
\textbf{题目}\par
在弱场横向无迹（TT）规范，沿$z$传播的线偏振引力波在探测器平面满足$h_{xx}(t)=h_0\cos\Omega t$、$h_{yy}(t)=-h_0\cos\Omega t$，$|h_0|\ll1$。Michelson探测器两臂沿$x,y$且未扰动长度均为$L$，探测激光真空波长$\lambda$；假设$\Omega L/c\ll1$。
\par\smallskip\noindent (1) 从线性化度规求两臂瞬时固有长度变化。
\par\smallskip\noindent (2) 用局部自由落体系测地偏离方程复核差动长度响应。
\par\smallskip\noindent (3) 求往返激光的主阶差动相位信号。
\par\smallskip\noindent (4) 说明长波近似、引力波偏振方向和高频响应函数的限制。
\par\medskip\hrule\medskip
\section*{M228\quad Thomas进动的Fermi–Walker输运与相对论轨道几何相位}
\textbf{题目}\par
自旋陀螺在平直闵可夫斯基时空中以实验室恒定速率$v<c$绕半径$R$的圆轨道运动，$\boldsymbol v(t)$切向、$\boldsymbol a(t)$径向向内，轨道角速率$\Omega=v/R$，假定陀螺无外加自旋转矩。
\par\smallskip\noindent (1) 用瞬时共动惯性系的连续非共线Lorentz推动说明自旋会有Thomas进动。
\par\smallskip\noindent (2) 写实验室Thomas进动角速度的通式并计算圆周运动的方向和大小。
\par\smallskip\noindent (3) 求完成一周实验室轨道后的总Thomas进动角。
\par\smallskip\noindent (4) 给出低速展开，说明其与普通动力学扭矩造成的进动区别。
\par\medskip\hrule\medskip
\chapter{电磁学}
\section*{E220\quad London涡旋的磁场、磁通量量子化与下临界场}
\textbf{题目}\par
第二类超导体在London极限$\kappa=\lambda_L/\xi\gg1$，考虑沿$z$轴的一条孤立直涡旋，穿透深度$\lambda_L$、芯半径$\xi$，通量子$\Phi_0=h/(2e)$，忽略其他涡旋及边界。
\par\smallskip\noindent (1) 写含点源的London磁场方程，求轴对称磁场分布。
\par\smallskip\noindent (2) 证明单涡旋携带的总磁通为$\Phi_0$并求绕向屏蔽电流。
\par\smallskip\noindent (3) 用场能和超流动能求涡旋每单位长度能量的对数主项。
\par\smallskip\noindent (4) 在外加场$H$下推导下临界场的主阶，并说明芯截止和相互作用限制。
\par\medskip\hrule\medskip
\section*{E221\quad 金属—介质界面表面等离激元的TM色散与光锥}
\textbf{题目}\par
平面界面$z=0$，上方$z>0$为无损介质$\epsilon_d>0$，下方$z<0$为各向同性金属$\epsilon_m(\omega)<0$；磁导率均为$\mu_0$。电场沿$xz$面，传播沿$x$，真空波数$k_0=\omega/c$。
\par\smallskip\noindent (1) 构造界面两侧指数衰减的TM模并写界面条件。
\par\smallskip\noindent (2) 解出表面波传播常数与局域衰减长度。
\par\smallskip\noindent (3) 对无耗散Drude金属$\epsilon_m(\omega)=1-\omega_p^2/\omega^2$，求长波与高波数极限。
\par\smallskip\noindent (4) 讨论TE模的缺失、能量局域和小损耗对波数的影响。
\par\medskip\hrule\medskip
\section*{E222\quad 圆柱导线的高频趋肤效应与Bessel复阻抗}
\textbf{题目}\par
半径$a$的均匀无限长导线，电导率$\sigma$、磁导率$\mu$，流经轴向谐波电流$I(t)=\operatorname{Re}(I_0e^{-i\omega t})$，忽略位移电流且材料线性。
\par\smallskip\noindent (1) 从Maxwell方程和Ohm定律推导圆柱内轴向电场微分方程。
\par\smallskip\noindent (2) 给出满足轴心正则性的Bessel解并求截面总电流。
\par\smallskip\noindent (3) 求每单位长度复阻抗的精确表达式。
\par\smallskip\noindent (4) 分析低频直流与高频强趋肤极限，给出平均焦耳损耗。
\par\medskip\hrule\medskip
\section*{E223\quad 双载流子Drude模型的Hall补偿与非饱和磁电阻}
\textbf{题目}\par
电子密度$n$、迁移率$\mu_e>0$和空穴密度$p$、迁移率$\mu_h>0$的均匀金属薄片，外磁场$B\hat z$，两载流子各服从独立弛豫时间Drude方程，$e>0$。
\par\smallskip\noindent (1) 叠加两类载流子求平面电导张量。
\par\smallskip\noindent (2) 求纵向和Hall电阻率。
\par\smallskip\noindent (3) 推导弱场Hall系数并说明Hall符号反转条件。
\par\smallskip\noindent (4) 对完全补偿$n=p$且$\mu_e=\mu_h$分析磁电阻随场的增长规律。
\par\medskip\hrule\medskip
\section*{E224\quad 带电枢电感的Faraday圆盘发电机与机电阻尼特征值}
\textbf{题目}\par
均匀磁场$B\hat z$中半径$R$的导体圆盘绕轴以角速度$\omega(t)$转动，轴与轮缘构成外接回路电阻$\mathcal R>0$、电感$L>0$。转子总转动惯量$J>0$，忽略摩擦和涡流，设$k=BR^2/2$。
\par\smallskip\noindent (1) 推出转盘轴到轮缘的动生电动势与电流方程。
\par\smallskip\noindent (2) 由磁力矩和角动量守恒求闭合机电耦合方程。
\par\smallskip\noindent (3) 消元得到转速的二阶常系数方程及过阻尼/欠阻尼阈值。
\par\smallskip\noindent (4) 证明机械能与磁储能之和严格按焦耳热率下降，并给无电感极限。
\par\medskip\hrule\medskip
\section*{E225\quad 双层二维电子气的对称与反对称等离激元色散}
\textbf{题目}\par
两平行二维电子层位于$z=\pm d/2$，层间距$d>0$，每层平衡面密度$n_s$、电子质量$m$、电荷$-e$，背景绝对介电率$\epsilon$，忽略惯性外的碰撞与量子压力，研究静电近似下平面波$\exp[i(qx-\omega t)]$，$q>0$。
\par\smallskip\noindent (1) 从电子连续性与动量方程求层面电荷密度对层电势的响应。
\par\smallskip\noindent (2) 从二维Coulomb Green函数构造两层电势耦合矩阵。
\par\smallskip\noindent (3) 分离同相、反相模式并求两支等离激元频率。
\par\smallskip\noindent (4) 求$qd\ll1$与$qd\gg1$的渐近形式，并解释反相声学支。
\par\medskip\hrule\medskip
\section*{E226\quad 短偶极天线的辐射角分布、总功率与辐射电阻}
\textbf{题目}\par
真空中沿$z$轴的短线性电偶极矩$\boldsymbol p(t)=p_0\cos\omega t\,\hat z$，$\omega>0$，$p_0$为实常数；远场$r\gg\lambda=2\pi c/\omega$，忽略材料损耗和高阶多极矩。
\par\smallskip\noindent (1) 从迟滞电磁势写出主导$1/r$的辐射电磁场。
\par\smallskip\noindent (2) 求时间平均的角向Poynting通量与方向图。
\par\smallskip\noindent (3) 对全立体角积分得到总辐射功率，并与Larmor公式交叉验证。
\par\smallskip\noindent (4) 对电流幅值$I_0$、电偶极有效长度$\ell$满足$p_0=I_0\ell/\omega$的短天线求辐射电阻。
\par\medskip\hrule\medskip
\section*{E227\quad 谐振腔微扰理论的广义本征值频移与介电体插入位置}
\textbf{题目}\par
无耗散、非色散、各向同性介质的封闭理想导体谐振腔满足$\nabla\times\mu^{-1}\nabla\times\boldsymbol E=\omega^2\epsilon(\boldsymbol r)\boldsymbol E$。假设模为非简并、$\mu$保持不变，在小体积$V_p$内介电常数增加$\delta\epsilon(\boldsymbol r)$且扰动足够弱。
\par\smallskip\noindent (1) 由Maxwell广义Hermitian本征值问题导出一阶频移公式。
\par\smallskip\noindent (2) 对一维理想腔驻波$E_0(x)=E_a\sin(\pi x/L)$求局部小介电片位置依赖频移。
\par\smallskip\noindent (3) 若位置从场节点移向反节点，分析符号、最大幅度与电场能量集中区。
\par\smallskip\noindent (4) 讨论简并模式、色散材料和强扰动为什么不能直接采用该公式。
\par\medskip\hrule\medskip
\section*{E228\quad 射频Paul离子阱的Mathieu参数稳定与赝势频率}
\textbf{题目}\par
电荷$q\ne0$、质量$m>0$的离子处于理想四极射频电场，电势$\Phi(x,y,z,t)=[U+V\cos(\Omega t)](x^2+y^2-2z^2)/(2r_0^2)$，$r_0,\Omega>0$。忽略离子间Coulomb作用、边缘场及射频耗散。
\par\smallskip\noindent (1) 验证电势满足无自由电荷区域的Laplace方程并列出三个坐标的运动方程。
\par\smallskip\noindent (2) 将运动化为标准Mathieu方程，确定其$a$、$q_M$参数（注意电荷$q$与Mathieu参数不同）。
\par\smallskip\noindent (3) 在$U=0$且小射频参数条件下推导三维赝势与长期世俗振动频率。
\par\smallskip\noindent (4) 说明静电四极势为何不能稳定三个方向，而射频阱如何绕过静态Earnshaw限制及近似失效情形。
\par\medskip\hrule\medskip
\chapter{热学}
\section*{T218\quad 一维Ising链的转移矩阵、磁化率与有限温度无序}
\textbf{题目}\par
$N$个周期性边界自旋$s_j=\pm1$，哈密顿量$\mathcal H=-J\sum_js_js_{j+1}-h\sum_js_j$，$J>0$，温度$T>0$，热力学极限$N\to\infty$。
\par\smallskip\noindent (1) 建立$2\times2$转移矩阵并求两个本征值。
\par\smallskip\noindent (2) 求每自旋自由能与精确磁化曲线。
\par\smallskip\noindent (3) 在$h=0$时求两点关联、关联长度和线性磁化率。
\par\smallskip\noindent (4) 说明有限温度无自发磁化及$T\to0$的奇异极限。
\par\medskip\hrule\medskip
\section*{T219\quad 热电发电器的焦耳分热与最大效率负载}
\textbf{题目}\par
均匀单级热电模块的Seebeck系数$\alpha$、内电阻$R$和零电流热导$K$均为常数，冷热端$T_h>T_c>0$，$\Delta T=T_h-T_c$；忽略Thomson效应和接触热阻，外接纯电阻$R_L$。
\par\smallskip\noindent (1) 从电压与负载匹配求电流、输出功率以及最大功率条件。
\par\smallskip\noindent (2) 结合Peltier热流与焦耳分热写热端吸热和制冷端放热，核对第一定律。
\par\smallskip\noindent (3) 以负载比$m=R_L/R$推导效率函数并给最大效率的最优$m$。
\par\smallskip\noindent (4) 写无量纲$ZT$形式，讨论$ZT\to0$与$ZT\to\infty$及可逆性。
\par\medskip\hrule\medskip
\section*{T220\quad Debye声子谱的低温热容、零点能与高温极限}
\textbf{题目}\par
三维各向同性晶体有$N$个原胞，声学声子在Debye近似中用截止频率$\omega_D$表征，三支声学模的总数为$3N$，定义$\Theta_D=\hbar\omega_D/k_B$。
\par\smallskip\noindent (1) 从三维态密度的$\omega^2$标度确定Debye态密度和截止归一化。
\par\smallskip\noindent (2) 计算声子热激发内能与零点能。
\par\smallskip\noindent (3) 推导$T\ll\Theta_D$时的$T^3$热容定律及其系数。
\par\smallskip\noindent (4) 检验$T\gg\Theta_D$时的Dulong–Petit极限并说明模型局限。
\par\medskip\hrule\medskip
\section*{T221\quad 有简并度二能级系统的Schottky热容峰与最优温标}
\textbf{题目}\par
$N$个彼此独立的粒子，各有基态能量0、简并度$g_0>0$与激发态能量$\Delta>0$、简并度$g_1>0$，温度$T$，无粒子数交换。
\par\smallskip\noindent (1) 求单粒子配分函数和激发态占有几率。
\par\smallskip\noindent (2) 推导热内能与定容热容的封闭表达。
\par\smallskip\noindent (3) 对$g_0=g_1$求热容峰位置的超越方程及唯一正根。
\par\smallskip\noindent (4) 分析低温和高温渐近，解释简并度对热容峰的影响。
\par\medskip\hrule\medskip
\section*{T222\quad 自旋半顺磁体的绝热退磁曲线、熵和Schottky热容}
\textbf{题目}\par
$N$个互不相互作用的磁矩大小$\mu>0$，能量为$-\mu Bs$，$s=\pm1$，外加磁场$B>0$、温度$T>0$；忽略晶格热容和磁相互作用。
\par\smallskip\noindent (1) 求配分函数、内能与平均磁化。
\par\smallskip\noindent (2) 计算绝对熵和定场热容。
\par\smallskip\noindent (3) 证明可逆绝热退磁满足$T/B$不变。
\par\smallskip\noindent (4) 讨论$B\to0$形式极限与第三定律、实际退磁制冷的限制。
\par\medskip\hrule\medskip
\section*{T223\quad 二维理想费米气体的常态密度与非解析低温化学势}
\textbf{题目}\par
面积$A$中二维非相对论自由费米子质量$m$、内禀简并度$g$，温度$T>0$，粒子总数$N$固定；忽略相互作用和边缘态，能量零点为单粒子色散最低值。
\par\smallskip\noindent (1) 由动量空间计数得到二维常量态密度$D$并定义$E_F$。
\par\smallskip\noindent (2) 精确求任意$T$的化学势$\mu(T)$。
\par\smallskip\noindent (3) 计算内能积分及低温热容主项。
\par\smallskip\noindent (4) 对比二维与三维费米气体的低温化学势变化并给高温经典极限。
\par\medskip\hrule\medskip
\section*{T224\quad Langevin–Ornstein–Uhlenbeck过程的涨落耗散与位移交叉极限}
\textbf{题目}\par
质量$m>0$的布朗粒子在温度$T$的均匀热库中作一维运动，动量阻尼系数$\gamma>0$，动力学满足$m\,dv=-\gamma v\,dt+\sqrt{2\gamma k_BT}\,dW_t$，其中$W_t$为标准Wiener过程，假设初速度已处于平衡分布。
\par\smallskip\noindent (1) 求平均速度的松弛时间和速度方差的平衡值。
\par\smallskip\noindent (2) 求平稳速度自关联$\langle v(t)v(0)\rangle$。
\par\smallskip\noindent (3) 积分得到任意时间的平均平方位移。
\par\smallskip\noindent (4) 求短时弹道和长时扩散极限，验证Einstein扩散系数。
\par\medskip\hrule\medskip
\section*{T225\quad 自引力孤立星团的Virial平衡与负微正则热容}
\textbf{题目}\par
$N\gg1$个质量$m$的经典粒子构成近似球形、自束缚且统计稳态的星团，忽略质量流失、辐射和外部压力，粒子间势能总和$U<0$为反平方引力对应的$-1$次齐次势，平均动能$K$。
\par\smallskip\noindent (1) 用Virial定理推导动能与引力势能的稳态关系。
\par\smallskip\noindent (2) 以速度分布温度定义$K=3Nk_BT/2$，求总能量及微正则热容。
\par\smallskip\noindent (3) 解释星团能量损失为何可使核心升温，并与普通正热容系统对照。
\par\smallskip\noindent (4) 讨论有限体积、外压、二体逃逸及正则系综对结论的限制。
\par\medskip\hrule\medskip
\chapter{振动与波动}
\section*{W226\quad 半无限弹性介质的Rayleigh表面波与本征方程}
\textbf{题目}\par
均匀各向同性线弹性体占$z<0$，自由表面$z=0$，密度$\rho$、Lamé常数$\lambda,\mu$，设$\mu>0$且$\lambda+2\mu>0$。沿$x$传播平面应变无耗散Rayleigh波，纵横波速$c_p=\sqrt{(\lambda+2\mu)/\rho}$、$c_s=\sqrt{\mu/\rho}$。
\par\smallskip\noindent (1) 用标量势和横向矢势构造向内部衰减的位移解。
\par\smallskip\noindent (2) 写自由表面剪切/法向应力为零的二元线性系统。
\par\smallskip\noindent (3) 导出Rayleigh特征方程并化为速度多项式。
\par\smallskip\noindent (4) 分析物理根范围、深度局域性和不可压极限。
\par\medskip\hrule\medskip
\section*{W227\quad 周期声学双层晶格的Bloch带隙与阻抗失配}
\textbf{题目}\par
一维理想无耗散声学超晶格由两种均匀介质逐层交替构成，厚度$d_1,d_2>0$，体密度$\rho_j$、声速$c_j$、平面波阻抗$Z_j=\rho_jc_j$，周期$a=d_1+d_2$，频率$\omega>0$。
\par\smallskip\noindent (1) 从声压和质点速度构造每层传播矩阵并说明行列式。
\par\smallskip\noindent (2) 求单周期Bloch色散的迹公式。
\par\smallskip\noindent (3) 用四分之一波厚度说明阻抗失配打开声学带隙。
\par\smallskip\noindent (4) 求禁带中有限周期数$N$的场包络衰减尺度并说明无失配极限。
\par\medskip\hrule\medskip
\section*{W228\quad 受张力Euler–Bernoulli梁的色散与自由端自然边界}
\textbf{题目}\par
均匀细梁弯曲刚度$EI>0$、质量线密度$\mu>0$，沿轴向有恒定拉力$T\ge0$，忽略阻尼和剪切形变。横向小位移$y(x,t)$。
\par\smallskip\noindent (1) 从动能、弯曲能和拉伸能变分推导梁方程。
\par\smallskip\noindent (2) 求无限长梁的弯曲—张力混合色散、相速与群速。
\par\smallskip\noindent (3) 给出两种回复机制交叉波数并分析长短波极限。
\par\smallskip\noindent (4) 对一端固支、另一端自由的有限梁，严格推导四个边界条件。
\par\medskip\hrule\medskip
\section*{W229\quad Sine–Gordon拓扑孤子、Lorentz收缩与精确能量}
\textbf{题目}\par
连续场$\phi(x,t)$满足$\phi_{tt}-c^2\phi_{xx}+\omega_0^2\sin\phi=0$，$c,\omega_0>0$，能量密度$\mathcal E=\phi_t^2/2+c^2\phi_x^2/2+\omega_0^2(1-\cos\phi)$。
\par\smallskip\noindent (1) 利用静态一阶积分求连接0与$2\pi$的拓扑kink。
\par\smallskip\noindent (2) 构造以恒定速度$|v|<c$传播的移动解。
\par\smallskip\noindent (3) 求静止kink能量和移动kink的Lorentz能量。
\par\smallskip\noindent (4) 证明拓扑荷守恒并说明该解不同于小振幅线性波。
\par\medskip\hrule\medskip
\section*{W230\quad 均匀Bragg光栅的耦合模方程、禁带与反射率}
\textbf{题目}\par
一维无耗散弱折射率周期调制使同频前行、后行波振幅$A(z),B(z)$满足耦合模方程$dA/dz=i\delta A+i\kappa B$、$dB/dz=-i\delta B-i\kappa A$，其中实失谐$\delta$、耦合强度$\kappa>0$，光栅长$L$，远端无后向入射$B(L)=0$。
\par\smallskip\noindent (1) 建立矩阵形式并求传播本征指数。
\par\smallskip\noindent (2) 推出禁带条件和场包络衰减常数。
\par\smallskip\noindent (3) 求有限光栅的反射率与在Bragg共振$\delta=0$的简式。
\par\smallskip\noindent (4) 检验短光栅、长光栅及离带条件下的能量守恒。
\par\medskip\hrule\medskip
\section*{W231\quad 旋转浅水方程的惯性重力波与零频地转模}
\textbf{题目}\par
$f$平面近似下，均匀层深$H>0$的浅水流体，重力加速度$g>0$、Coriolis参数$f\ne0$。线性水平速度$(u,v)$和自由面扰动$\eta$满足$u_t-fv=-g\eta_x$、$v_t+fu=-g\eta_y$、$\eta_t+H(u_x+v_y)=0$。
\par\smallskip\noindent (1) 将方程写成平面波本征值问题。
\par\smallskip\noindent (2) 求两支惯性重力波色散及零频模。
\par\smallskip\noindent (3) 从线性潜在涡度守恒说明地转模的物理来源。
\par\smallskip\noindent (4) 求群速度及长短波极限，并区别地转模式与传播波。
\par\medskip\hrule\medskip
\section*{W232\quad 聚焦非线性Schrödinger平面波的调制不稳定性}
\textbf{题目}\par
一维无损包络满足$i\psi_t+P\psi_{xx}+Q|\psi|^2\psi=0$，$P,Q$为实数且$P\ne0$；考虑均匀背景振幅$A_0>0$。
\par\smallskip\noindent (1) 求恒定幅值的平面波背景及相位频率。
\par\smallskip\noindent (2) 以实虚扰动线性化并导出侧带色散关系。
\par\smallskip\noindent (3) 给出调制不稳定的参数区与最大增长率、最易增长波数。
\par\smallskip\noindent (4) 讨论散焦介质、守恒量及模型小扰动适用性。
\par\medskip\hrule\medskip
\section*{W233\quad 时变张力弦的Mathieu参数激发与阻尼阈值}
\textbf{题目}\par
一根两端固定、长度$L$、线密度$\rho_\ell$的理想弦，张力受小周期调制$T(t)=T_0[1+h\cos(\Omega t)]$，$|h|\ll1$，并允许每模具有弱线性阻尼$2\beta_n\dot q_n$。
\par\smallskip\noindent (1) 展开固定端本征空间模态并得到各模Mathieu方程。
\par\smallskip\noindent (2) 在主参数共振$\Omega\simeq2\omega_n$附近用慢变幅法求增长率。
\par\smallskip\noindent (3) 求频率失谐导致的不稳定舌边界。
\par\smallskip\noindent (4) 加入阻尼求主共振最低驱动阈值并讨论能量来源。
\par\medskip\hrule\medskip
\chapter{光学}
\section*{O226\quad 平面—凹面激光谐振腔的ABCD稳定性与高斯模}
\textbf{题目}\par
一条轴向长度$L$的空腔，一端为平面理想反射镜，另一端为曲率半径$R>0$的凹面反射镜，腔内折射率1，波长$\lambda$，近轴标量、高反射率且无介质损耗。
\par\smallskip\noindent (1) 建立往返光线矩阵并从特征值求稳定条件。
\par\smallskip\noindent (2) 求基模腰位置、瑞利长度和束腰半径。
\par\smallskip\noindent (3) 从横向Gouy相位推导模频率及纵模简并结构。
\par\smallskip\noindent (4) 检验共焦边界、近乎平面和同心式退化边界的光学意义。
\par\medskip\hrule\medskip
\section*{O227\quad 环形孔径的Fraunhofer衍射、中心遮挡与极限零点}
\textbf{题目}\par
振幅均匀照明的环形理想孔径，外半径$a$，内半径$b=\epsilon a$，$0\le\epsilon<1$，真空波长$\lambda$，远场观察角$\theta$；设$u=ka\sin\theta$、$k=2\pi/\lambda$。
\par\smallskip\noindent (1) 从圆对称傅里叶—Bessel变换推导环孔复振幅。
\par\smallskip\noindent (2) 以轴上场作归一化，求强度及总开孔面积变化。
\par\smallskip\noindent (3) 求第一暗纹所满足的超越方程，并说明无中心遮挡极限。
\par\smallskip\noindent (4) 推导窄环极限的归一化场及其第一暗环与普通圆孔的差异。
\par\medskip\hrule\medskip
\section*{O228\quad 正弦相位光栅的Bessel衍射级、能量守恒与零级抑制}
\textbf{题目}\par
单色相干平面光正入射无损薄相位光栅，周期$d$，复透过率$t(x)=\exp[i m\cos(2\pi x/d)]$，相位调制深度$m\ge0$，忽略传播中的偏振和材料吸收。
\par\smallskip\noindent (1) 用Jacobi–Anger展开推导离散衍射级复振幅。
\par\smallskip\noindent (2) 给出各级次相对能量及传播角条件。
\par\smallskip\noindent (3) 在弱调制极限推导零级与$\pm1$级强度分配并检验能量守恒。
\par\smallskip\noindent (4) 求零级被完全抑制的第一个调制深度并讨论更高级次能量。
\par\medskip\hrule\medskip
\section*{O229\quad 单总线微环共振器的临界耦合与吸收线宽}
\textbf{题目}\par
单总线侧耦合光学微环采用时间耦合模理论，腔本征角频率$\omega_0$，内损耗振幅衰减率$\gamma_i\ge0$、总线外耦合振幅衰减率$\gamma_e>0$，输入输出波振幅满足功率归一化$|s|^2=P$。
\par\smallskip\noindent (1) 建立共振腔振幅方程及输入—输出关系。
\par\smallskip\noindent (2) 求通带幅透射系数与功率透射谱。
\par\smallskip\noindent (3) 推导临界耦合条件与总功率去向。
\par\smallskip\noindent (4) 求谱线全宽、加载品质因数，并讨论无内损极限。
\par\medskip\hrule\medskip
\section*{O230\quad 弱导引阶跃光纤的LP模截止与单模V数}
\textbf{题目}\par
芯半径$a$的圆柱阶跃折射率光纤，芯折射率$n_1>n_2>0$，外包层无限，真空波长$\lambda$；采用弱导引标量近似$\Delta=(n_1-n_2)/n_1\ll1$。
\par\smallskip\noindent (1) 从横向Helmholtz方程给出芯内、包层径向Bessel函数解。
\par\smallskip\noindent (2) 用界面连续条件导出LP模式超越本征方程。
\par\smallskip\noindent (3) 推导归一化频率$V$以及单模条件$V<2.405$的来源。
\par\smallskip\noindent (4) 讨论传播常数范围、模式截止和近似适用性。
\par\medskip\hrule\medskip
\section*{O231\quad 单光子干涉和NOON态的量子Fisher信息与相位极限}
\textbf{题目}\par
两路径构成理想无损相位干涉仪，相位$\phi$仅作用于第二路径。考虑单光子态$|\psi_1(\phi)\rangle=(|1,0\rangle+e^{i\phi}|0,1\rangle)/\sqrt2$，以及两模式N光子NOON态$|\psi_N(\phi)\rangle=(|N,0\rangle+e^{iN\phi}|0,N\rangle)/\sqrt2$。
\par\smallskip\noindent (1) 从纯态QFI公式求单光子相位信息量。
\par\smallskip\noindent (2) 求$N$个独立单光子的总QFI和标准量子极限。
\par\smallskip\noindent (3) 求NOON态的QFI与Heisenberg标度。
\par\smallskip\noindent (4) 在每个光子独立透过率$\eta<1$的后选择探测下求有效信息的损失因子。
\par\medskip\hrule\medskip
\section*{O232\quad 二次谐波的相位失配、sinc效率与准相位匹配}
\textbf{题目}\par
无泵浦耗尽、慢变包络与平面波近似下，二次谐波振幅$A_2(z)$满足$dA_2/dz=i\kappa A_1^2 e^{i\Delta k z}$，$A_1$为常量实泵浦包络、$\kappa$实数，$\Delta k=2k_1-k_2$。
\par\smallskip\noindent (1) 积分求长度$L$晶体的二次谐波复幅。
\par\smallskip\noindent (2) 推导相位失配对生成强度的sinc平方规律和相干长度。
\par\smallskip\noindent (3) 在非线性系数每个相干长度变号的准相位匹配中求首谐波有效耦合。
\par\smallskip\noindent (4) 讨论泵浦耗尽、Manley–Rowe光子数关系与理想近似的界限。
\par\medskip\hrule\medskip
\section*{O233\quad 全反射界面的Goos–Hänchen位移与反射相位色散}
\textbf{题目}\par
TE偏振平面光由非磁介质$n_1$射向较低折射率介质$n_2$，$n_1>n_2>0$，入射角$\theta>\arcsin(n_2/n_1)$，真空波数$k_0$。研究有限宽度准单色光束的反射光斑横向位移。
\par\smallskip\noindent (1) 写出全反射区的切向波数、法向波数和倏逝常数。
\par\smallskip\noindent (2) 求TE Fresnel反射复系数及其相位。
\par\smallskip\noindent (3) 用驻相法推导Goos–Hänchen横向位移。
\par\smallskip\noindent (4) 讨论临界角与掠射极限的形式发散及有限束宽修正。
\par\medskip\hrule\medskip
\chapter{近代物理}
\section*{Q226\quad 对称双吸引delta势的偶奇束缚态与隧穿劈裂}
\textbf{题目}\par
一维粒子质量$m$，$V(x)=-g[\delta(x-a)+\delta(x+a)]$，$g>0,a>0$，$\hbar$保留。束缚态能量$E=-\hbar^2\kappa^2/(2m)$，$\kappa>0$。
\par\smallskip\noindent (1) 由波函数连续和导数跳跃推导偶态、奇态的束缚能方程。
\par\smallskip\noindent (2) 证明奇态出现需要满足最小势阱间距条件，偶态始终存在。
\par\smallskip\noindent (3) 在$a\kappa_0\gg1$时求两态能量的指数劈裂，$\kappa_0=mg/\hbar^2$。
\par\smallskip\noindent (4) 讨论$a\to0$时两delta融合的单势阱极限及态数变化。
\par\medskip\hrule\medskip
\section*{Q227\quad 二维质量Dirac带的Berry曲率与整数量子化正则化}
\textbf{题目}\par
考虑二维连续两带哈密顿量$H(\boldsymbol k)=v(k_x\sigma_x+k_y\sigma_y)+M(k)\sigma_z$，$v>0$；先取常质量$M(k)=m\ne0$，再用$M(k)=m-B|\boldsymbol k|^2$作高动量正则化，$B\ne0$。以Berry联络$\mathcal A=i\langle u_-|\nabla_{\boldsymbol k}u_-\rangle$定义下能带曲率。
\par\smallskip\noindent (1) 求两能带及常质量下能带的Berry曲率分布。
\par\smallskip\noindent (2) 在连续动量平面积分曲率，解释为什么仅得半整数通量。
\par\smallskip\noindent (3) 对Wilson型$M=m-Bk^2$计算紧化平面的Chern数并指出相变。
\par\smallskip\noindent (4) 讨论Berry曲率与Hall电导符号、规范选择及晶格正则化之间的关系。
\par\medskip\hrule\medskip
\section*{Q228\quad 四次微扰量子振子的能级非等间距与选择定则}
\textbf{题目}\par
一维谐振子质量$m$、角频率$\omega$，添加弱正四次势$V_1=\lambda x^4$，$\lambda>0$，假设在所研究量子数$n$范围内一阶微扰有效。
\par\smallskip\noindent (1) 用升降算符建立$x^4$的矩阵元结构与选择规则。
\par\smallskip\noindent (2) 求任意量子数$n$的一阶能级修正。
\par\smallskip\noindent (3) 求相邻能级跃迁频率随$n$的变化。
\par\smallskip\noindent (4) 给出高激发态下微扰适用的小参数并检验$\lambda\to0$极限。
\par\medskip\hrule\medskip
\section*{Q229\quad Breit–Wigner单共振散射与Wigner时间延迟}
\textbf{题目}\par
单弹性散射通道的孤立共振能量$E_0$、总宽度$\Gamma>0$，背景相移在共振附近可视为常数$\delta_{\rm bg}$。定义单位模散射矩阵$S(E)=e^{2i\delta_{\rm bg}}[E-E_0-i\Gamma/2]/[E-E_0+i\Gamma/2]$。
\par\smallskip\noindent (1) 证明弹性单通道散射矩阵的幺正性并求共振相位。
\par\smallskip\noindent (2) 由能量导数推导Wigner–Smith延迟的Lorentz线形和峰值。
\par\smallskip\noindent (3) 对两端口共振透射给出峰值上界及对称耦合条件。
\par\smallskip\noindent (4) 比较$\Gamma\to0$的时间—能量尺度与背景相位的影响。
\par\medskip\hrule\medskip
\section*{Q230\quad 投影量子Zeno效应、非选择测量与强退相干极限}
\textbf{题目}\par
二能级体系哈密顿量$H=\hbar\Omega\sigma_x/2$，初态为$\sigma_z$的$-1$本征态$|g\rangle$，总演化时间$t$，每经过$\delta t=t/N$进行一次理想$\{|g\rangle,|e\rangle\}$投影测量。
\par\smallskip\noindent (1) 求每段自由演化后仍在基态的几率。
\par\smallskip\noindent (2) 求$N$次均得到$g$的选择性存活概率及其$N\to\infty$极限。
\par\smallskip\noindent (3) 对不做选择、只记录最终布居的重复测量给出基态概率。
\par\smallskip\noindent (4) 用Markov型强退相干Lindblad方程求有效慢跃迁率并讨论与投影极限的联系。
\par\medskip\hrule\medskip
\section*{Q231\quad 高频驱动二能级Floquet系统与Bessel抑制的有效隧穿}
\textbf{题目}\par
周期性驱动二能级哈密顿量$H(t)=\frac\Delta2\sigma_x+A\cos(\omega t)\sigma_z$，$\Delta>0$、$\omega>0$，并假设$\hbar\omega\gg\Delta$且可控制驱动幅度$A$。
\par\smallskip\noindent (1) 构造消除纵向驱动的周期性酉变换。
\par\smallskip\noindent (2) 求变换后横向耦合的谐波展开。
\par\smallskip\noindent (3) 在高频平均近似下推导有效Floquet哈密顿量。
\par\smallskip\noindent (4) 给出相干隧穿抑制的Bessel零点条件及近似适用边界。
\par\medskip\hrule\medskip
\section*{Q232\quad 时变谐振子的Ermakov–Lewis不变量与非平衡量子压缩}
\textbf{题目}\par
单位质量经典振子满足$\ddot x+\omega(t)^2x=0$，$\omega(t)>0$；辅助尺度$\rho(t)>0$满足Ermakov方程$\ddot\rho+\omega(t)^2\rho=\rho^{-3}$。在量子情形取$[x,p]=i\hbar$且$p=\dot x$形式对应哈密顿量$H(t)=p^2/2+\omega(t)^2x^2/2$。
\par\smallskip\noindent (1) 构造经典Ermakov–Lewis不变量并证明严格守恒。
\par\smallskip\noindent (2) 用新时标把变量频率振子映射为单位频率谐振子。
\par\smallskip\noindent (3) 写出对称排序的量子不变量及其本征值结构。
\par\smallskip\noindent (4) 在频率恒定与突变调制极限核对不变量的物理意义。
\par\medskip\hrule\medskip
\section*{Q233\quad Jaynes–Cummings腔量子电动力学的真空Rabi劈裂与失谐转移}
\textbf{题目}\par
单模量子腔频率$\omega_c$，二能级原子频率$\omega_q$，旋转波近似哈密顿量$H=\hbar\omega_ca^\dagger a+\hbar\omega_q\sigma_z/2+\hbar g(a\sigma_++a^\dagger\sigma_-)$，$g>0$。
\par\smallskip\noindent (1) 证明总激发数守恒并写每个固定激发子空间矩阵。
\par\smallskip\noindent (2) 求 dressed state 能级和共振真空Rabi劈裂。
\par\smallskip\noindent (3) 求初始原子激发态$|e,n\rangle$转移到$|g,n+1\rangle$的时间依赖几率。
\par\smallskip\noindent (4) 求大失谐色散频移并说明旋转波与二能级近似限制。
\par\medskip\hrule\medskip
\end{document}
